% Successor-seeded scalar iteration: a self-contained arithmetic bootstrap.
% Requires amsmath, amssymb, amsthm and theorem/lemma/proposition environments.
% No document wrapper. All induction invariants below are literal equations.

\section{A successor-seeded scalar basis with two computation schemas}
\label{js:section}

This section completes the explicitly unresolved one-variable
successor-seeded candidate in the supplied shared-continuation
report~\cite{shared}, in its section ``A genuinely smaller numerical
grammar: the precise unresolved step.''  The candidate basis and its
conditional initializer compiler are inherited from that report.
Unlike the counter-aware successor recursor, its update has just one
changing input.  The initializer remains a separate argument, but the
iterator starts at its successor.  The missing result is a numerical
predecessor term with a whole-domain proof, obtained here by adapting
the supplied scalar $2/3$ factor-map construction to shifted arithmetic
and correcting a saturated predecessor at three small arguments.

\subsection{Syntax, computations, and the precise resources used}

The terms are finite staged expressions
\[
 t ::= x\mid 0\mid J(n,a;u.h).
\]
Only $h$ binds $u$; its other free variables are passive parameters.
Define the abbreviation
\[
 Sa:=J(0,a;u.0),\qquad 1:=S0,
\]
and use standard numerals as finite strings of this abbreviation.
There are exactly two numerical computation schemas:
\begin{align}
 J(0,a;u.h)&=Sa,\tag{J0}\label{js:j0}\\
 J(Sn,a;u.h)&=h[J(n,a;u.h)/u].\tag{JS}\label{js:js}
\end{align}
The first schema, with the abbreviation expanded, says that
$J(0,a;u.h)=J(0,a;u.0)$.  It does not presuppose a separate successor
constructor.  No function variables or uniform evaluator are present.

The bootstrap uses equality replacement, numerical cases
\[
 \frac{\Gamma\vdash F(0)\qquad\Gamma\vdash F(Sz)}
      {\Gamma\vdash F(t)},
 \tag{NC}\label{js:nc}
\]
and contextual induction only on a literal equation $E$:
\[
 \frac{\Gamma\vdash E(0)\qquad\Gamma,E(z)\vdash E(Sz)}
      {\Gamma\vdash E(t)}.
 \tag{EI}\label{js:ei}
\]
The ordinary eigenvariable conditions apply.  In an ordinary
intuitionistic first-order presentation, NC follows from coverage
$\forall x\,(x=0\vee\exists y\,x=Sy)$.  In the native paired-binder
presentation it is the already available case rule.  No use of
successor injectivity, predecessor, pairing, an equality characteristic,
or induction on a compound formula is made before predecessor is
constructed below.

The finite range splits used below require the distinction proved in
Section~\ref{nb:section}. Functional range elimination follows from NC;
contextual finite cases with retained induction hypotheses use the finite
cancellation and guard construction of Lemma~\ref{nb:guarded-cases}.
The cycle identities are established first by functional range elimination
alone. Thus a displayed finite range is not an uncharged disjunction or
an assumption of the ordinary-connective macros.

We shall use the following admissible consequence of EI.  If two fixed
updates are pointwise equal, and the displayed initial values agree,
their iterates agree at every counter.  Induct on the literal equation
between the two iterates: J0 supplies its base and JS supplies its step.
This is a derived uniqueness argument, not a primitive extensionality
rule for update bodies.  All passive arguments are kept fixed while
performing the induction.

\subsection{A partial selector and shifted addition}

Put
\[
 B(n,a,b):=J(n,a;u.b),\qquad Z(n):=B(n,0,0),
 \qquad a\oplus b:=J(b,a;u.Su).
\]
Their computation endpoints are
\[
 B(0,a,b)=Sa,\quad B(Sn,a,b)=b,\qquad
 Z(0)=1,\quad Z(Sn)=0,
\]
\[
 a\oplus0=Sa,\qquad a\oplus Sb=S(a\oplus b).
\]
The operation $\oplus$ has standard meaning $a+b+1$.  Ordinary
unshifted addition has not been defined and is not used in the proof.
Literal-equation induction proves
\[
 0\oplus a=Sa,\quad Sa\oplus b=S(a\oplus b),\quad
 a\oplus b=b\oplus a,\quad
 (a\oplus b)\oplus c=a\oplus(b\oplus c).
 \tag{Shifted algebra}\label{js:oplus}
\]
For the first two equations induct on the second argument.  For
commutativity induct on $b$, using the first two equations.  For
associativity induct on $c$; at zero both sides are
$S(a\oplus b)$, and the successor step applies $S$ to the previous
equality.  These proofs do not cancel a successor.

In particular, the following are genuine available terms:
\[
 D(a):=S^2(a\oplus a),\qquad
 T(a):=S^4((a\oplus a)\oplus a),\qquad
 R(a):=S^3(a\oplus a).
 \tag{Affine terms}\label{js:affine}
\]
Their standard meanings are respectively $2a+3$, $3a+6$, and
$2a+4$.  The proof uses their displayed term definitions.  From
\eqref{js:oplus} one obtains
\[
 D(0)=3,\quad D(Sa)=S^2D(a),\qquad
 T(0)=6,\quad T(Sa)=S^3T(a).
 \tag{Affine shifts}\label{js:affine-shifts}
\]

\subsection{Two finite cycles}\label{js:cycles}

Define
\[
 N(n):=Z\bigl(J(n,0;u.Z(u))\bigr).
\]
The inner iterator is always zero or one: at zero it is one, and at a
successor it is a value of $Z$.  The endpoints of $Z$ and finite cases
therefore give
\[
 N(0)=0,\quad N(Sn)=Z(N(n)),\quad N(S^2n)=N(n),
 \qquad N(n)\in\{0,1\}.
 \tag{Parity cycle}\label{js:parity}
\]
This asserts a proved two-state recurrence, not a pre-existing parity
operation.  For example, if $K(n)=J(n,0;u.Z(u))$, then
$N(Sn)=Z(Z(K(n)))=K(n)=Z(N(n))$ by its two cases.

Next put
\[
 j(u):=B\bigl(u,0,B(Z(N(u)),1,0)\bigr),\qquad
 M(n):=j\bigl(j(J(n,0;u.j(u)))\bigr).
\]
Finite computations give
\[
 j(0)=1,\qquad j(1)=2,\qquad j(2)=0.
\]
The range of $j$ is contained in $\{0,1,2\}$: its zero branch is
one, and its successor branch is zero or two according to $N(u)$.
Consequently the definitions and these three computations give
\[
 M(0)=0,\quad M(Sn)=j(M(n)),\quad M(S^3n)=M(n),
 \qquad M(n)\in\{0,1,2\}.
 \tag{Three-cycle}\label{js:three-cycle}
\]
Here $M(0)=j(j(1))=0$, and the successor equation follows by
associativity of composing the fixed term $j$ with itself.  The period
three follows from its three possible values.  No numerical division
algorithm has been assumed.

\subsection{The shifted ramp and factor map}

Set
\[
 G(n):=J(n,0;u.u\oplus M(S^2u)).
\]
Its first values are
\[
 G(0)=1,\ G(1)=2,\ G(2)=4,\ G(3)=5,\
 G(4)=7,\ G(5)=8,\ G(6)=10.
\]

\begin{lemma}[Ramp identities]\label{js:ramp}
The preceding theory proves
\[
 M(G(n))=S(N(n)),\qquad G(S^2n)=S^3G(n),
 \tag{Ramp}\label{js:ramp-identities}
\]
and admits the following two finite-case calculations:
\[
 \begin{array}{c|cc}
 &N(n)=0&N(n)=1\\ \hline
 G(Sn)&S G(n)&S^2G(n)\\
 G(S^3n)&S^4G(n)&S^5G(n).
 \end{array}
 \tag{Ramp cases}\label{js:ramp-cases}
\]
\end{lemma}
\begin{proof}
Induct on the single equation $M(G(n))=S(N(n))$.
Its base is $M(1)=1$.  If $N(n)=0$, the induction equation gives
$M(G(n))=1$, hence $M(S^2G(n))=0$ by the three-cycle.  The next
state is $G(n)\oplus0=S G(n)$, whose $M$-value is two, as
required.  If $N(n)=1$, the current $M$-value is two and
$M(S^2G(n))=1$.  The next state is
$G(n)\oplus1=S^2G(n)$, whose $M$-value is one.  This proves
the invariant and the first line of the case table.  Applying that line
twice, with the parity cycle, gives the two-step equation; applying it
three times gives the last line.
\end{proof}

\begin{lemma}[Affine identities]\label{js:affine-lemmas}
Throughout the entire numerical domain,
\[
 N(D(a))=1,\qquad N(T(a))=N(a),\qquad
 S G(D(a))=T(a),
 \tag{Affine parity and ramp}\label{js:affine-parity}
\]
\[
 G(T(a))=T(G(a))\oplus N(a).
 \tag{Ramp scaling}\label{js:ramp-scaling}
\]
\end{lemma}
\begin{proof}
For $N(D(a))=1$, use EI, $D(0)=3$, and the two-step parity
period.  For $N(T(a))=N(a)$ use EI, $T(0)=6$, and
$T(Sa)=S^3T(a)$; three parity steps toggle both sides once.

For $S G(D(a))=T(a)$ use EI on that literal equation.  The base
is $S G(3)=6$.  The successor step uses
$D(Sa)=S^2D(a)$ and $G(S^2v)=S^3G(v)$, so it applies three
successors to the previous equation, just as $T$ does.

For \eqref{js:ramp-scaling}, the base is
$G(6)=10=T(1)\oplus0$.  Retain its equation as the induction
hypothesis, and split $N(a)$.  In the zero case, $N(T(a))=0$,
so the new left side is $S^4G(T(a))$.  The old right side is
$T(G(a))\oplus0=S T(G(a))$, while the new right side is
\[
 T(SG(a))\oplus1=S^5T(G(a)).
\]
Thus both sides have changed by four successors.  In the one case,
the new left side is $S^5G(T(a))$.  The old right side is
$T(G(a))\oplus1=S^2T(G(a))$, while the new right side is
\[
 T(S^2G(a))\oplus0=S^7T(G(a)).
\]
Thus both sides have changed by five successors.  These calculations
use only \eqref{js:oplus}, \eqref{js:affine-shifts}, and the ramp
case table; no successor has been cancelled.
\end{proof}

Define the factor map
\[
 H(u):=B(Z(N(u)),G(u),u).
\]
Its two cases are
\[
 N(u)=0\ \Longrightarrow\ H(u)=u,\qquad
 N(u)=1\ \Longrightarrow\ H(u)=S G(u).
\]

\paragraph{Why these shifts are useful.}
For intuition on standard natural numbers only, $G(u)=\lfloor3u/2\rfloor+1$.
Thus $H$ fixes even inputs and sends an odd $u$ to $3(u+1)/2$.
After shifting the state by three, an active step multiplies it by $3/2$.
The affine maps $D(a)=2a+3$ and $T(a)=3a+6$ are precisely doubling
and tripling in this shifted coordinate: $D(a)+3=2(a+3)$ and
$T(a)+3=3(a+3)$. Starting at two makes the unshifted base even and
fixed by $H$, while its shifted value is five. The later seed sequence
therefore has standard values $5\cdot2^x-3$. This motivates a comparator
that tests when its powers of two have been exhausted. The proof below
uses the displayed $J$ terms and literal equations; division, exponentiation,
and this standard-value description are not prerequisites.

\begin{lemma}[The factor-map equations]\label{js:factor}
The following global equations hold:
\[
 H(2)=2,\qquad H(D(a))=T(a),\qquad H(T(a))=T(H(a)).
 \tag{Factor}\label{js:factor-equations}
\]
\end{lemma}
\begin{proof}
The first equation follows from $N(2)=0$.  Since $N(D(a))=1$,
the second is $S G(D(a))=T(a)$.
For the last equation split $N(a)$, noting that $N(T(a))=N(a)$.
In the zero case both sides equal $T(a)$.  In the one case,
\eqref{js:ramp-scaling} gives
\[
 H(T(a))=S G(T(a))
 =S\bigl(T(G(a))\oplus1\bigr)
 =S^3T(G(a))=T(SG(a))=T(H(a)).
\]
\end{proof}

\subsection{A comparator without stripping an initializer}

Put
\[
 A(x):=J(x,0;u.R(u)),\qquad E(x):=S A(x),
\]
\[
 F(x,y):=J(y,A(x);u.H(u)),\qquad
 L(x,y):=Z(N(F(x,y))).
\]
By direct computation and shifted algebra,
\[
 E(0)=2,\qquad E(Sx)=D(E(x)),\qquad F(x,0)=E(x).
\]
The definition of $F$ is the decisive initializer alignment:
its successor-seeded initial value is exactly $E(x)=S A(x)$.
There is no hidden use of predecessor to initialize it at $E(x)$.

\begin{lemma}[Comparator equations]\label{js:comparator}
The term $L$ is zero or one and satisfies
\[
 L(0,y)=1,\qquad L(Sx,0)=0,\qquad
 L(Sx,Sy)=L(x,y).
 \tag{Comparator}\label{js:comparator-equations}
\]
\end{lemma}
\begin{proof}
Its range follows from the range of $Z$.  EI on $y$ proves
$F(0,y)=2$, since $E(0)=2$ and $H(2)=2$; hence $L(0,y)=1$.
Also $F(Sx,0)=E(Sx)=D(E(x))$ has parity one, giving
$L(Sx,0)=0$.

For the remaining equation first prove, by EI on $y$ with $x$ fixed,
\[
 F(Sx,Sy)=T(F(x,y)).
\]
At zero, the left side is
$H(E(Sx))=H(D(E(x)))=T(E(x))$.  In the successor step both
iterates apply $H$, and $H\circ T=T\circ H$ rewrites the result
into $T(F(x,Sy))$.  Applying $Z\circ N$ and using
$N(T(v))=N(v)$ gives the desired comparator equation.

This proof avoids cancellation in an equation between two successor
initializers.  Such cancellation would be unavailable at this stage.
\end{proof}

\subsection{Saturation and a finite correction give true predecessor}

Define conditional increment by
\[
 K(c,u):=B(c,u,u),\qquad
 B_m(u):=K(Z(L(Su,m)),u),\qquad
 V(t,m):=J(t,0;u.B_m(u)).
\]
Thus $K(0,u)=Su$ and $K(Sv,u)=u$.  The counter $m$ is a passive
parameter of the iteration defining $V$.

\begin{lemma}[Shifted saturation]\label{js:saturation}
The theory proves
\[
 V(St,S^2k)=S V(t,Sk),\qquad
 V(n,Sn)=Sn,\qquad
 V(Sn,Sn)=V(S^2n,Sn)=Sn.
 \tag{Saturation}\label{js:saturation-equations}
\]
\end{lemma}
\begin{proof}
Numerical cases on $c$ give $K(c,Su)=S K(c,u)$.  Consequently
the comparator equations give
\[
 B_{Sm}(Su)=S B_m(u),\qquad B_{S^2k}(1)=2.
\]
EI on $t$, with $k$ fixed, now proves the first displayed equation.
At zero its two sides equal two.  At a successor, apply the defining
update to the induction equation and then use the conjugacy
$B_{S^2k}(Su)=S B_{Sk}(u)$.

EI gives $L(Sn,n)=0$ directly from the comparator equations.
A separate EI proves $V(n,Sn)=Sn$: the base is $V(0,1)=1$,
and the successor step uses the previously proved uniform shift at
$t=k=n$.  This is substitution into a completed lemma, not a change
of a passive parameter inside an active induction hypothesis.
Finally
\[
 B_{Sn}(Sn)=K(Z(L(S^2n,Sn)),Sn)
           =K(Z(L(Sn,n)),Sn)=K(1,Sn)=Sn.
\]
One and then two more iterator steps give the last two equations.
\end{proof}

Put
\[
 Q(m):=J\bigl(m,0;u.K(Z(L(S^2u,m)),u)\bigr).
\]
This term is a capped predecessor rather than the desired predecessor:
\[
 Q(0)=Q(1)=1,\qquad Q(S^2k)=Sk.
 \tag{Capped predecessor}\label{js:capped}
\]
The first two equations are direct computations.  For the last,
\[
 L(S^2u,S^2k)=L(u,k)=L(Su,Sk)
\]
identifies its update with $B_{Sk}(u)$.  The derived iterator uniqueness
lemma, proved at an independent counter $t$ with $k$ fixed, gives
\[
 J(t,0;u.K(Z(L(S^2u,S^2k)),u))=V(t,Sk).
\]
Substitute $t=S^2k$ and apply \eqref{js:saturation-equations}.

Now define one fixed correction term
\[
 \boxed{\quad
 P(m):=B\bigl(Z(L(3,m)),\ Q(Q(m)),\ L(2,m)\bigr).
 \quad}\tag{Predecessor}\label{js:predecessor}
\]
There is no arbitrary selector hidden in this expression.  Its outer
operation is the original partial selector $B$, whose zero branch adds
exactly one successor.

\begin{theorem}[Predecessor and initializer recovery]\label{js:recovery}
The two-schema successor-seeded scalar theory proves
\[
 P(0)=0,\qquad P(Sn)=n.
\]
For every fixed update $h$, the finite term
\[
 I(n,a;v.h):=
 P\bigl(J(n,a;u.S(h[Pu/v]))\bigr)
 \tag{Initializer compiler}\label{js:initializer}
\]
satisfies
\[
 I(0,a;v.h)=a,\qquad I(Sn,a;v.h)=h[I(n,a;v.h)/v].
\]
\end{theorem}
\begin{proof}
The comparator equations give the three direct values
\[
 P(0)=0,\qquad P(1)=0,\qquad P(2)=1.
\]
Indeed, at these inputs $L(3,m)=0$, so the outer $B$ selects
$L(2,m)$, whose three values are zero, zero, one.
At $m=S^3k$, one has $L(3,m)=1$.  Thus the zero-control branch
of $B$, together with \eqref{js:capped}, gives
\[
 P(S^3k)=S Q(Q(S^3k))=S Q(S^2k)=S^2k.
\]
Numerical cases twice on $n$ now give $P(Sn)=n$ from these displayed
equations.  Alternatively, two literal-equation inductions, first on
$P(S^2n)=Sn$ and then on $P(Sn)=n$, give the same conclusion.
No earlier cancellation or injectivity was used.  Successor injectivity
can now be derived by applying $P$ to an equality of successors.

For the initializer compiler, J0 and the proved predecessor equation
give the base $P(Sa)=a$.  At a successor, JS exposes the term
$S(h[P(J(n,a;u.S(h[Pu/v])))/v])$.  The outer $P$ removes its
displayed successor and leaves exactly $h[I(n,a;v.h)/v]$.
There is no invariant about the internal state to prove: the displayed
successor in the update is sufficient.
\end{proof}

\subsection{Consequences and the bi-interpretation boundary}

The construction supplies the full initialized scalar basis $0,S,I$.
In particular the unrestricted selector is now the derived term
$C(n,a,b):=I(n,a;u.b)$, with its ordinary endpoints $a$ and $b$.
The supplied scalar arithmetic development can therefore be used from
this point onward to obtain actual-equality reflection, induction on
arbitrary quantifier-free formulas, pairing and projections, and the
compiler for every fixed general primitive-recursion description.
Those later arguments are not premises of this bootstrap.

In ordinary iPRA, each fixed $J$ description is interpreted by primitive
recursion with seed $Sa$ and its displayed scalar update.  Both
computation schemas, coverage, separation, and literal-equation induction
are then provable.  Conversely \eqref{js:initializer}, followed by the
existing scalar compiler, interprets every fixed iPRA operation.
For each fixed description, literal-equation induction proves uniqueness
of its defining recurrence after earlier subdescriptions have been
identified.  Applied to both composite translations, this identifies
each numerical operation with its double translation throughout the
entire original domain.  Both domain formulas are $x=x$ and both
numerical isomorphism graphs are $x=y$.

For the ordinary presentations, full intuitionistic logic is preserved.
For the native paired-binder presentation, its ordinary logical macros
and the formula return-translation proof must still be included; they
apply once the derived selector $C$ is available.  This establishes
bi-interpretability only when that logical bridge and the scalar
arithmetic completion accompany the present argument.  It does not
replace them by agreement merely on the standard natural numbers.

The arithmetic advance is completion of the predecessor obligation
left open for this exact candidate in the shared-continuation
report~\cite{shared}: the affine factor-map alignment, capped
predecessor, and finite correction \eqref{js:predecessor} supply its
missing proof.  The candidate grammar and initializer compiler are not
new.  Finite cycles, a $2/3$ factor map, and saturation are adapted from
the supplied ordinary scalar predecessor construction.  No global
literature priority is asserted.

The now-completed basis has two primitive numerical constructors and
two computation schemas, with a single update variable.  A smaller
complete proof checker remains a separate implementation claim, which
must charge the actual decoder, scope checks, logical rules, and driver.
